\documentclass[11pt]{article}

\usepackage[a4paper,margin=28mm]{geometry}
\usepackage{amsmath,amssymb}
\usepackage{booktabs}

\title{Diagonal Coordinate Normalization for the Bundle-Adjustment Schur System}
\author{}
\date{}

\begin{document}
\maketitle

\section{Key insight}

The camera, point, rotation, translation, focal-length, and distortion
coordinates in bundle adjustment can have very different numerical scales.
Consequently, the diagonal entries of the normal matrix may span many orders of
magnitude.  Solving that system directly makes an approximate inner method
sensitive to floating-point reduction order, pseudoinverse thresholds, and a
stopping criterion measured in poorly scaled coordinates.

The useful operation is a change of coordinates.  Let the diagonal matrices
$D_c$ and $D_p$ contain square roots of the camera and point normal-matrix
diagonals.  We solve for dimensionless variables $z_c,z_p$ and recover the
physical increments by
\begin{equation}
  \delta_c = D_c^{-1} z_c,
  \qquad
  \delta_p = D_p^{-1} z_p.
\end{equation}
The transformed normal matrix has approximately unit diagonal.  This is useful
for both PCG and Nesterov; it is not specific to either inner algorithm.

\section{Original linearized problem}

For residual vector $r$, camera Jacobian $J_c$, and point Jacobian $J_p$, a
Gauss--Newton step minimizes
\begin{equation}
  \min_{\delta_c,\delta_p}
  \left\|r + J_c\delta_c + J_p\delta_p\right\|_2^2.
\end{equation}
Define
\begin{equation}
  U = J_c^{\mathsf T}J_c,
  \qquad
  V = J_p^{\mathsf T}J_p,
  \qquad
  W = J_c^{\mathsf T}J_p,
\end{equation}
and
\begin{equation}
  b_c = -J_c^{\mathsf T}r,
  \qquad
  b_p = -J_p^{\mathsf T}r.
\end{equation}
The normal equations are
\begin{equation}
  \begin{bmatrix}
    U & W \\
    W^{\mathsf T} & V
  \end{bmatrix}
  \begin{bmatrix}
    \delta_c \\
    \delta_p
  \end{bmatrix}
  =
  \begin{bmatrix}
    b_c \\
    b_p
  \end{bmatrix}.
\end{equation}
Eliminating the points gives
\begin{align}
  S\delta_c &= b_S, \\
  S &= U-WV^{-1}W^{\mathsf T}, \\
  b_S &= b_c-WV^{-1}b_p.
\end{align}
After solving for $\delta_c$, point back-substitution is
\begin{equation}
  \delta_p = V^{-1}\left(b_p-W^{\mathsf T}\delta_c\right).
\end{equation}

\section{Diagonal coordinate transform}

Choose a small floor $\tau>0$ and define
\begin{align}
  D_c &= \operatorname{diag}\left(
    \max\left(\sqrt{\left|\operatorname{diag}(U)\right|},\tau\right)
  \right), \\
  D_p &= \operatorname{diag}\left(
    \max\left(\sqrt{\left|\operatorname{diag}(V)\right|},\tau\right)
  \right).
\end{align}
In the implementation, $\tau=10^{-10}$.  Let
\begin{equation}
  T = \operatorname{blkdiag}(D_c,D_p),
  \qquad
  \begin{bmatrix}\delta_c\\\delta_p\end{bmatrix}
  = T^{-1}
  \begin{bmatrix}z_c\\z_p\end{bmatrix}.
\end{equation}
Substitution into the quadratic model and multiplication by $T^{-1}$ produce
the symmetric congruence transform
\begin{equation}
  H' = T^{-1}HT^{-1},
  \qquad
  b' = T^{-1}b.
\end{equation}
Its blocks are
\begin{align}
  U' &= D_c^{-1}UD_c^{-1}, \\
  V' &= D_p^{-1}VD_p^{-1}, \\
  W' &= D_c^{-1}WD_p^{-1}, \\
  b_c' &= D_c^{-1}b_c, \\
  b_p' &= D_p^{-1}b_p.
\end{align}
The transformed Schur equations are therefore
\begin{align}
  S'z_c &= b_S', \\
  S' &= U'-W'(V')^{-1}(W')^{\mathsf T}, \\
  b_S' &= b_c'-W'(V')^{-1}b_p'.
\end{align}
Point back-substitution and conversion to physical coordinates are
\begin{align}
  z_p &= (V')^{-1}\left(b_p'-(W')^{\mathsf T}z_c\right), \\
  \delta_c &= D_c^{-1}z_c, \\
  \delta_p &= D_p^{-1}z_p.
\end{align}
For an exact undamped solve in exact arithmetic, this transformation does not
change the physical Gauss--Newton step.  It only changes the basis in which the
linear system is represented.

\section{Damping in normalized coordinates}

The implementation first normalizes the Gauss--Newton blocks and then applies
local damping:
\begin{align}
  U'_{\lambda} &= (1+\lambda)U' + \lambda\epsilon I, \\
  V'_{\lambda} &= (1+\lambda)V' + \lambda\epsilon I.
\end{align}
The identity term is thus isotropic in dimensionless coordinates.  When mapped
back to the original variables, it corresponds to the diagonal metric induced
by $D_c$ and $D_p$, rather than an identity metric over quantities with
incompatible physical units.  The nonlinear candidate cost is still evaluated
in the original camera and point coordinates.

For PALM acceptance, the corresponding block-diagonal damping penalty is
evaluated directly as
\begin{equation}
  \lambda\left(
    \|J_c\delta_c\|_2^2 + \|J_p\delta_p\|_2^2
    + \epsilon\|z_c\|_2^2 + \epsilon\|z_p\|_2^2
  \right).
\end{equation}
Computing the two Jacobian-vector norms directly avoids relying on a copied
sparse normal matrix and makes the evaluated penalty identical to the metric
used by the normalized damping equations.

This detail means that the normalized damped step need not equal the old
unnormalized damped step.  That difference is intentional: damping becomes
comparable across coordinate families instead of being dominated by the
largest raw scales.

\section{Why approximate solvers improve}

Although an exact undamped solution is invariant, a finite inner iteration is
not invariant to coordinate scaling.  Normalization improves the practical
solve for several reasons:

\begin{enumerate}
  \item The diagonal of $U'$ and $V'$ is approximately one, reducing the dynamic
        range seen by sparse matrix products.
  \item Relative residual and step stopping tests become meaningful because one
        coordinate family no longer dominates the Euclidean norm.
  \item Batched pseudoinverses are less sensitive to scale-dependent singular
        value thresholds.
  \item Roundoff and nondeterministic GPU sparse-reduction order are less likely
        to be amplified by a severely ill-scaled Schur complement.
  \item A fixed Nesterov step scale is much more appropriate for the normalized
        operator.
\end{enumerate}

PCG still uses the camera block inverse as its Krylov preconditioner, and the
Nesterov iteration still uses the camera block inverse in its gradient map.
Those are algorithmic preconditioners.  The $D_c,D_p$ operation described here
is different: it changes the coordinates of the entire camera--point system,
including $U$, $V$, $W$, both right-hand sides, back-substitution, and final step
recovery.

\section{Observed effect}

For Ladybug \texttt{problem-1064-113655-pre}, five overlap-refined partitions,
20 epochs, and the original PALM acceptance condition, the observed results
were
\begin{center}
\begin{tabular}{lrr}
  \toprule
  Inner method & Final objective & Time (s) \\
  \midrule
  PCG, unnormalized & approximately $1.9\times 10^6$ to $2.9\times 10^6$ & $25$--$33$ \\
  PCG, normalized & $587050.84$ & $10.96$ \\
  Nesterov, normalized & $586578.39$ & $12.23$ \\
  CPU reference & approximately $591006$ & --- \\
  \bottomrule
\end{tabular}
\end{center}
The normalized runs generally required only ten inner iterations per block.
Repeated normalized PCG runs also agreed to nearly machine precision, whereas
the unnormalized trajectory was sensitive to small GPU numerical differences.

\section{Implementation summary}

For each local Schur solve:
\begin{enumerate}
  \item Assemble the original $U$, $V$, and $W$ blocks.
    \item Compute $D_c^{-1}$ and $D_p^{-1}$ from the block diagonals.
    \item Apply the two-sided scaling to $U$, $V$, and $W$.
    \item Scale $b_c$ and $b_p$ once on the left.
    \item Apply damping and solve the transformed Schur system with PCG or
        Nesterov.
    \item Back-substitute for $z_p$ in transformed coordinates.
    \item Recover $\delta_c$ and $\delta_p$ with $D_c^{-1}$ and $D_p^{-1}$.
    \item Evaluate the damping metric from $J_c\delta_c$, $J_p\delta_p$, $z_c$,
      and $z_p$.
    \item Evaluate the candidate and acceptance condition in original
        coordinates.
\end{enumerate}

In \texttt{palm\_ba.py}, BAE local solves enable this transformation by
default.  The option \texttt{--no-inner-jacobi} disables it for controlled
comparisons.

\end{document}
